\begin{statementbox}
	\textbf{\textcolor{CStatement}{条件}}
	\begin{description}[style=nextline, leftmargin=2.8em, labelsep=0.5em, itemsep=0.35em]
		\item[\textbf{\textcolor{CStatement}{条件 1（定义域与运算）：}}]
			$f(x),g(x)$的定义域均关于原点对称；运算在公共定义域上进行；除法运算要求分母不为零。
	\end{description}
	\textbf{\textcolor{CStatement}{结论}}
	\begin{description}[style=nextline, leftmargin=2.8em, labelsep=0.5em, itemsep=0.45em]
		\item[\textbf{\textcolor{CStatement}{结论一（加减规则）：}}]
			\textbf{\textcolor{CStatement}{直观描述：}}
			奇函数之间相加减仍为奇函数，偶函数之间相加减仍为偶函数。
			\[
				\begin{aligned}
					&f(-x)=-f(x),\; g(-x)=-g(x)\;\Rightarrow\;(f\pm g)(-x)=-(f\pm g)(x),\\&f(-x)=f(x),\; g(-x)=g(x)\;\Rightarrow\;(f\pm g)(-x)=(f\pm g)(x).
			\end{aligned}
			\]
		\item[\textbf{\textcolor{CStatement}{结论二（乘法规则）：}}]
			\textbf{\textcolor{CStatement}{直观描述：}}
			奇函数相乘得偶函数，奇与偶相乘得奇函数，偶与偶相乘得偶函数。
			\[
				\begin{aligned}
					&f(-x)=-f(x),\; g(-x)=-g(x)\;\Rightarrow\;(fg)(-x)\\
					&\qquad=f(x)g(x)\\
					&\qquad=(fg)(x),\
					\[
						2mm]&f(-x)=-f(x),\; g(-x)=g(x)\;\Rightarrow\;(fg)(-x)\\
						&\qquad=-f(x)g(x)\\
						&\qquad=-(fg)(x),\
						\[
							2mm]&f(-x)=f(x),\; g(-x)=g(x)\;\Rightarrow\;(fg)(-x)\\
							&\qquad=f(x)g(x)\\
							&\qquad=(fg)(x).
					\end{aligned}
					\]
				\item[\textbf{\textcolor{CStatement}{推广结论(除法规则):}}]
					\textbf{\textcolor{CStatement}{直观描述:}}
					除法与乘法规则完全一致,但需确保分母恒不为零.
					\[
						\begin{aligned}
							&f(-x)=-f(x),\; g(-x)=-g(x),\; g(x)\neq0\;\Rightarrow\;\left(\frac{f}{g}\right)(-x)\\
							&\qquad=\frac{f(x)}{g(x)}\\
							&\qquad=\left(\frac{f}{g}\right)(x),\
							\[
								2mm]&f(-x)=-f(x),\; g(-x)=g(x),\; g(x)\neq0\;\Rightarrow\;\left(\frac{f}{g}\right)(-x)\\
								&\qquad=-\frac{f(x)}{g(x)}\\
								&\qquad=-\left(\frac{f}{g}\right)(x),\
								\[
									2mm]&f(-x)=f(x),\; g(-x)=g(x),\; g(x)\neq0\;\Rightarrow\;\left(\frac{f}{g}\right)(-x)\\
									&\qquad=\frac{f(x)}{g(x)}\\
									&\qquad=\left(\frac{f}{g}\right)(x).
							\end{aligned}
							\]
					\end{description}

				\end{statementbox}
